% !TeX root = ../papers/nc-p3.tex
\section{交代・対称な印付き窓と可換な長さ四の環}
\label{sec:mixed-point-frobenius-algebras}
有限点枝の印付き表示で $g=q=3$ のとき、
定理\ref{thm:rank-eight-finite-symmetry}の交換子の証明のうち
反対階数8を使うのは混合型を除外する段のみである。以下ではこの混合型そのものを分類する。
$U,V$ を三次元の標準内積空間とし、$V$ に外積を選ぶ。
記号の向きを固定して $G=F$ を交代側、$Q=P$ を対称側に置く。逆の配置では開始位置を移す。
\begin{theorem}\label{thm:mixed-point-frobenius-algebras}
混合型の印付きデータは、四次元可換単位的代数 $K$ と線型汎関数 $\tau$ で
$\tau(1)=1$, $\tau(ab)$ が完全であるもの、および $U=\ker\tau$ 上の
相似までの非退化対称形式 $m$ に対応する。印を保つ添字保持同型の必要十分条件は
$\tau$ と $[m]$ を保つ単位代数同型の存在である。
$U$ の内積を $x\cdot z=-\tau(xz)$ とし、
$m(x,z)=(Lx)\cdot(Lz)$ となる同型 $L:U\to V$ を選ぶ。
各局所因子のsocleの生成元 $d\in K$ について
\[
      y_d=(L^t)^{-1}\big(d-\tau(d)1\big)\in V
\]
を作る。元の窓が標準Koszul左右AS正則であることと、全ての $d$ について
$y_d\cdot y_d\ne0$ であることは同値である。局所因子は非被約でもよい。
特に四辺が有限な混合型候補は全て正則で、反対階数10、最小周期2である。
\end{theorem}

\begin{proof}[長さ四の環と二周期表示]
$G_a=[La]_\times$ と置く。三独立な交代行列を持つので $L$ は可逆である。
$n_b=L^t b$ と置けば印付きの交換子条件は
\[
 [Q_b,Q_d]=n_b n_d^t-n_d n_b^t
\]
となる。Jacobi恒等式から $Q_b n_d-Q_d n_b=h_b n_d-h_d n_b$ となる線型形式 $h$ が一意にある。
実際、$b,d$ を添字対 $23,31,12$ として左辺を列に持つ行列を $J$、$n_i$ を列に持つ行列を $N$ とする。
上の交換子をJacobi恒等式に代入すれば $JN^t+NJ^t=0$ である。
三行の交代行列 $N^{-1}J$ はちょうど表示した $h$ で一意に書ける。
$B_{n_b}=Q_b-h_b I$ とおくと $B_z x=B_x z$ であり、$K=k1\oplus U$ に積
\[
 (r,x)(s,z)=(rs-x\cdot z,\ rz+sx+B_z x)
\]
を入れられる。各 $z$ の乗法行列
$\bigl(\begin{smallmatrix}0&-z^t\\ z&B_z\end{smallmatrix}\bigr)$ は互いに交換する。
非対角ブロックは上の対称性、右下は交換子式そのものである。
単位的可換な積で左乗法が互いに交換すれば結合律が従う。
スカラー部 $\tau$ の積対合は $1\oplus(-I_U)$ で完全である。
逆にこのような $(K,\tau)$ から $B_z$ は $U$ 上の乗法の $U$ 成分として復元され、
$Q_b$ は $B_{n_b}$ の跡零部分である。従って元の交換子式も成立する。

位置を偶・奇に折り畳むには、二位置先の偶の補文字 $r$ の符号を反転し、奇文字はそのまま用いる。
さらに奇の補座標を $s'=s+h_y$ と書き直す。隣接写像の六成分は、二つの順で
\[
           r y\mp (Lx)\times y,\qquad s'x+B_{n_y} x
\]
である。実際、最初の三成分は印付き四テンソルの $c=*,d=j$、残りは $c=i,d=*$ の列である。
反対向きの窓に移ると $G^t=-G$、$Q^t=Q$ が二つの符号を与える。
従って関係にはこの二組の行空間を取ればよい。
印付き同型公式で $F=G,P=Q$ の各三基底を保つには
$A^tA=B^t B=sI$ が必要十分である（ここだけ $B$ は奇位置の基底変換を表す）。
$A=a O,\ B=a P$、$O,P$ を直交変換とすれば
$L'=a\det(P)P L O^t$, $Q'_y=a O Q_{P^ty}O^t$ である。
従って $U$ 上の $\det(P)O$ が単位を固定して $(K,\tau)$ を対応させ、
$m$ は倍率 $a^2$ まで保たれる。逆にそのような代数同型を $O$ とし、
$L'=a P L O^t$ と書く。平方根 $a$ の符号で $\det P=1$ とできる。
$Q$ は上記の復元で対応するため、同じ印付き変換を得る。
\end{proof}

\begin{proof}[二次双対の構成と必要性]
$M=k\oplus V$ を積
$(r,u)(s,v)=(rs-u\cdot v,rv+su-u\times v)$ の分裂四元数環とする。
標数零の代数閉体上で $M\simeq M_2(k)$ であり、スカラー部は半跡、純成分の行列式は $v\cdot v$ である。
$D=M\oplus K$ で $\ell(m,k)=\operatorname{tr}(m)/2-\tau(k)$,
$H=\ker\ell$, $z=(1,1)$ と置く。点座標の二つの部分空間を
\[
 E=\{(r+Lx,x):r\in k,x\in U\},\qquad
 F=\{(y,s'+n_y):y\in V,s'\in k\}
\]
とする。$\ell(ab)$ は非退化で $F=E^\perp$。$E F=F E=H$（積の線型包）である。
積のスカラー部は両環で $-(Lx)\cdot y$、残る六成分は先の二周期写像である。
六成分の独立性は $r,s'$ の係数から、スカラー部の独立性も同じ二係数とその非零性から従う。
従って $H/kz$ への積が二次関係の評価そのものである。

どの一次元イデアルも一つの局所因子のsocleであり、その生成元 $d$ は $\tau(d)\ne0$ である。
これは有限Frobenius環の完全対合と $dK=kd$ による。$F$ の $K$ 射影は同型で、
$d$ の持上げ $f$ の行列成分は $y_d$ である。$y_d$ が非可逆なら $fD$ は次元3以下で
（行列側が高々2、可換側が1）、$\ell$ は $fD$ で非零である。
ゆえに $fE\subset fD\cap H$ の次元は2以下で、二次評価の階数も2以下となる。
片側対称表示に戻して定理\ref{thm:single-opposite-evaluation-ranks}を適用すれば非正則である。
また幾何学的にはこの評価に射影直線以上のファイバーがある。
\end{proof}

\begin{proof}[非被約な環を含む十分性]
以下全てのsocle持上げが可逆行列成分を持つとする。
$D$ の次数別コピーに $E,F$ を交互に掛けた部分からなる添字付き代数を $C$ とする。
各一次空間は $D$ の単元を稠密に含む。
$E$ の $K$ 成分 $U$ はどの極大イデアルとも異なる。もし一致すれば $\tau$ はそのイデアルを消し、
完全積対合に反する。$F$ では $K$ への射影が同型、行列成分の三平面にも単元がある。
$H$ の右安定化単元はスカラーのみである（完全対合で $H^\perp=k1$）。
従って $H$ に各次位置の一次空間の独立な二単元を掛けた二超平面は異なる。
長さ3の積は $D$ であり、以後も $D$ である。

$K$ のイデアル $I$ と $k^2$ の部分空間 $N$ に対し、$J_{I,N}$ を
$C$ の正次数成分と右イデアル $\{(X,u):\operatorname{im}X\subset N,\ u\in I\}$ の次数別交叉とする。
両始点で同じ記号を使う。二次元イデアル $I$ を任意に取り、$I'=\operatorname{ann}_K I$ とする。
$\dim I'=2$, $\operatorname{ann} I'=I$ は完全対合から従う。$\tau$ はどの非零イデアル上でも非零である。
$E,F$ のうち $K$ 成分を $I$ に制限した部分はともに二次元で、行列成分への射影は単射である。
$E$ 側は単位行列を含み、$F$ 側は純四元数の二平面で非退化三変数二次形式の全等方平面ではない。
従って行列式が零の射影根 $b$ があり、その行列 $X$ は階数1である。
この $b$ の $K$ 成分は $I$ を生成する。実際、生成しなければ積イデアルは零か一次元である。
$E$ 側の零成分は単位行列の倍、一次元イデアルは $U$ と交わらない。
$F$ 側は $K$ 射影が単射で、一次元イデアルでは仮定により行列が可逆である。
この議論は二重根と $I=I'$ も含む。

次位置の $I'$ に制限した行列二平面は、$I$ に制限した元位置の二平面の行列跡直交補である。
実際、$I I'=0$ と $F=E^\perp$、両二平面の次元による。
このうち像が $\ker X$ にある行列は一次元であり、そこに次の根 $c$ がある。
実際、今の二平面で $X$ と独立な可逆 $X'$ を取る。$\operatorname{im}Y\subset\ker X$ の二平面では
$\operatorname{tr}(XY)=0$ は自動、$\operatorname{tr}(X'Y)=0$ は非零の一式である。
従って $b$ の右零化空間は次の一次空間で $kc$ であり、二次以後では $D$ の右零化イデアルとの交叉である。
$b$ と次位置の一次空間の積は三次元、従って $J_{I,\operatorname{im}X}$ の二次成分全部である。
また $bH=bD$ である。その右零化イデアルに $I'$ の元で $\tau$ が非零になるものがあるからである。
一次でもこの行列像条件の解は $kb$ である（同じ二平面に可逆 $X'$ がある）。
ゆえに $bC=J_{I,\operatorname{im}X}$、その右零化イデアルは次位置の $cC$ となる。

ここからの線型イデアルの構成を記す。$I$ に制限した二平面で $X'$ に対応する元 $b'$ を加える。
$bC$ に対する $b'$ の右コロンは $J_{K,N}$, $N=(X')^{-1}\operatorname{im}X$ である。
このイデアルの次元は長さ1,2,3以後で $2,5,6$ である。一次では $E\to M$ は同型、
$F\to M$ は核一次元で像が跡零三平面、$\operatorname{im}Y\subset N$ の行列のうち跡零のものは一次元だからである。
従って短完全列で $(b,b')C=J_{I,k^2}$ が従う。右辺の次元も長さ1,2,3以後で $2,5,6$ であり、
$bC$ とコロン商の次元を足すとこれに達する。

この $J_{K,N}$ も一次生成元の追加で得られる。$X,X'$ の $N$ への制限は比例するので、
像が $N$ に入る行列二平面では二つの跡条件は一つの非零条件である。
従って次位置に $K$ 成分が $I'$ を生成し、行列像が $N$ の根 $c'$ がある。
同じ行列像条件の二平面で、$K$ 成分が単元の独立な $c''$ を取れる。
$F$ 側では $(0,1)$ が使える。$E$ 側では各剰余体への写像は、
反対位置の該当socleの持上げとの行列跡対合の非零定数倍である。
実際、$\ell(EF)=0$ と $\tau(kd)=\tau(d)\bar k$ である。
その行列は可逆なので、像が $N$ に入る行列全体ではこの対合は恒等零でない。
有限個の極大イデアルを避けられることが従う。
$c'C$ に対する $c''$ の右コロンは $J_{I',k^2}$ であり、同じ次元比較で
$(c',c'')C=J_{K,N}$ を得る。

最後に $I$ を二次元から一次元ずつ三次元のイデアル、さらに $K$ へ増やす。
各包含 $I_j\subset I_{j+1}$ に対し、$K$ 成分が $I_{j+1}/I_j$ で非零の一次生成元を追加する。
$U\cap I_j$ は $I_j$ の余次元1であり両始点でこの選択ができる。
右コロンはこの単純商を消す極大イデアル $\mathfrak m$ に対する $J_{\mathfrak m,k^2}$ である。
その商は各次数一次元であり、得られるイデアルは $J_{I_{j+1},k^2}$ に等しい。
全ての三次元イデアルが二次元イデアルを含むことは組成列から従う。
以上のイデアルを両始点で一緒に取れば、零から各イデアルへ一つずつ線型生成元を足す際の
コロンも同じ集合に属する。従って定理\ref{thm:half-period-semisimple}と同じ短完全列の
二重帰納で $C$ はKoszulである。この証明で $K$ の根基を割る必要はない。

$z$ は $C$ の両始点の二次に属する中心非零因子である。
従って $C/(z)$ はKoszul、各始点のHilbert関数は $1,4,6,4,1$ である。
最上次数の汎関数 $D/H\to k$ は $\ell$ で、$E,F$ の完全直交と $H^\perp=kz$ により
全ての相補次数で積対合は完全である。
二次の積はちょうど所与の六成分なので二次双対性で元の窓の標準Koszul左右AS正則性を得る。
\end{proof}

\begin{proof}[有限枝の階数と周期]
混合型の残余対合は定理\ref{thm:rank-eight-finite-symmetry}の式で非零なので階数は少なくとも9、
片側対称なので高々10である。階数9の場合でも全点で中支持4を仮定する必要はない。
この混合型の対称核は階数4であることを直接確かめる。
$D_G=-\operatorname{cof}(L)$ として、残余対合が階数1なら、
$I,Q_b$ と跡直交する対称行列の二平面中に、$D_G^tD_G$ の非零倍 $Y$ がある。
実際、残余対合の核は $D_G^{-t}Y D_G^{-1}$ がスカラーとなる条件そのものである。
従って $Y$ は $(L^tL)^{-1}$ の定数倍であり、後者に規格化する。
偶文字対の純成分に $Y$ を縮約すると $G$ 側の寄与は $-2I$ である
（外積積公式）。$Q$ 側は補題\ref{lem:ternary-net-commutator}から
$D_Q P_Q(Y)D_Q^t=(\det D_Q)Y$ で、$D_Q=-D_G^t=\operatorname{cof}(L)^t$ を入れると
$P_Q(Y)=I$ である。星印一個の成分は $\operatorname{tr}(YQ_b)=0$ で消え、
両星印側の $\operatorname{tr}Y$ も零である。ゆえに対称核の生成元はブロック対角 $1\oplus Y$ で階数4である。
有限枝では定理\ref{thm:smooth-symmetric-kernel-all}と両対称AS階数二分法により階数9は不可能である。
残る階数10で周期1なら偶位置の交換を一文字ずらして奇位置の交換も得るので両対称となり再び矛盾する。
既に二周期表示があるため最小周期は2である。
有限枝で印を忘れた同型は、相手の高々二十個の局所点のうち対応する混合型の単点を順に印として
上の $(K,\tau,[m])$ を比較すれば必要十分である。
\end{proof}
